\setcounter{section}{1}

\Needspace{5\baselineskip}
\hypertarget{gansux8badux7ec3ux7684ux6a21ux5f0fux5d29ux6e83ux95eeux9898ux4e0eux89e3ux51b3ux65b9ux6848}{%
\section{GANs训练的模式崩溃问题与解决方案}\label{gansux8badux7ec3ux7684ux6a21ux5f0fux5d29ux6e83ux95eeux9898ux4e0eux89e3ux51b3ux65b9ux6848}}

\Needspace{5\baselineskip}
\hypertarget{ux4e00ux6a21ux5f0fux5d29ux6e83}{%
\subsection{一、模式崩溃}\label{ux4e00ux6a21ux5f0fux5d29ux6e83}}

假设我们正在训练一个 GAN 来生成手写数字（如 MNIST 数据集，包含 0 到 9）。一个成功的 G（生成器）应该能生成所有 10 种不同的、逼真的数字（0, 1, 2, 3, 4, 5, 6, 7, 8, 9）。

如果 \(G\) 和 \(D\) ``对抗实力相同''，并``共同进化''，是不会发生这种情况的。但事实上判别器 \(D\) 在训练早期常学得太快，很快就识别出了 \(G\) 生成的是假的。由于 \(G\) 在各个方向生成得都不好，\(G\) 收到的梯度方向非常混乱，且前面讲述过，\(D\) 函数内部的 Sigmoid 函数会导致 \(p \to 0\) 时梯度消失，这些梯度的幅度比较小。

这时 \(G\) 可能在训练中偶然发现，自己生成的``7''这个数字特别逼真（概率显著大于 \(0\)），每次都能轻易地骗过 \(D\)（判别器）。由于 \(G\) 的唯一目标就是骗过 \(D\)，它会在这一片混乱的梯度中，生成``7''的方向作为唯一稳定的梯度方向，且往往幅度较大。这会把模型引向生成更多``7''的方向，这时便会有更多正反馈，更多梯度让模型进一步导向生成``7''。无论给它什么随机噪声 \(z\) 作为输入，它的输出都是一个（或一小组）看起来差不多的``7''。在这种情况下，我们就说 \(G\) ``崩溃''到了一个单一的模式（数字 ``7''）上，而完全忽略了数据集中的其他所有``模式''（0, 1, 2, 3\ldots）。

\Needspace{5\baselineskip}
\hypertarget{ux4e8cux6a21ux5f0fux5d29ux6e83ux7684ux89e3ux51b3ux65b9ux6848}{%
\subsection{二、模式崩溃的解决方案}\label{ux4e8cux6a21ux5f0fux5d29ux6e83ux7684ux89e3ux51b3ux65b9ux6848}}

\Needspace{5\baselineskip}
\hypertarget{ux4e00wgan}{%
\subsubsection{（一）WGAN}\label{ux4e00wgan}}

WGAN (Wasserstein GAN)，原始 GAN 的``对抗''太激烈了，尤其是判别器 \(D\) 变得太好时，会让生成器 \(G\) 无法学习。为了解决这个问题，WGAN 做了几个关键的改变。

\Needspace{4\baselineskip}
\begin{enumerate}
\def\labelenumi{\arabic{enumi}.}
\tightlist
\item
  改变D的工作和目标函数
\end{enumerate}

原始 GAN 中，\(D\) 是一个分类器 (Classifier)。输出一个 0 到 1 之间的概率（``这是假的'' vs ``这是真的''）；WGAN 中，\(D\) 被重新定义为一个评论家 (Critic)。它不再输出一个概率，而是输出一个分数 (Score)，撤掉了原来的 Sigmoid 函数。这个分数可以（理论上）是任何实数（比如 \(-10, 0, 50\)）。\(D\) 的目标变成了：给真实数据 \(x\) 尽量打高分，给伪造数据 \(G(z)\) 尽量打低分。\(G\) 的目标也随之改变：努力生成 \(G(z)\)，让 \(D\) 给它打出尽可能高的分。

\Needspace{4\baselineskip}
\begin{enumerate}
\def\labelenumi{\arabic{enumi}.}
\setcounter{enumi}{1}
\tightlist
\item
  梯度正则化
\end{enumerate}

\(D\) 的目标变成了给真实数据 \(x\) 尽量打高分，给伪造数据 \(G(z)\) 尽量打低分，但如果一味鼓励拉开高低分，会让 \(D\) 为了拉开高低分而让自己的打分曲线变得``过于陡峭''，导致梯度爆炸。原始方法是权重裁剪 (Weight Clipping)，这是 WGAN 原始论文中提出的方法，强行检查它的每一个权重，如果某个权重超过了一个小范围（比如 \(0.01\)），就把它强行``拉回''到 \(0.01\)；如果低于 \(-0.01\)，就把它``拉回''到 \(-0.01\)。这会导致逐层梯度相乘过程中发生梯度消失 (Vanishing Gradients)，且表达能力受限 (Limited Capacity)：\(D\) 为了满足这个约束，会``偷懒''，把大部分权重都推到边界值（\(0.01\) 或 \(-0.01\)），变成了一个非常简单的函数，只能用``好''或``坏''等词来打分，失去了学习复杂、精细打分曲线的能力。

\Needspace{5\baselineskip}
于是我们选择用如下数学公式作为惩罚项：

\[
\lambda \cdot (\lVert \nabla D(x) \rVert - 1)^2
\]

\Needspace{5\baselineskip}
这就迫使 \(D\) 在训练中必须时时刻刻调整自己的打分曲线，使其在任何地方的``坡度''都尽量保持在 \(1\) 左右。考虑到实际计算开销，我们一般选取真实与伪造的中间点计算惩罚项：

\begin{enumerate}
\def\labelenumi{\arabic{enumi}.}
\item
  取样：我们从数据集中随机取一个真实图像 \(x_{real}\)。
\item
  生成：我们让 \(G\) 生成一个伪造图像 \(x_{fake}\)。
\item
  插值：我们在这两个图像之间随机选择一个``中间点'' \(\hat{x}\)（读作 x-hat）。

  在数学上，我们随机选一个 0 到 1 之间的小数 \(\epsilon\)，然后计算 \(\hat{x}=\epsilon\cdot x_{real}+(1-\epsilon)\cdot x_{fake}\)。这就像在两个点的连线上随机找一个点。
\item
\Needspace{5\baselineskip}
  检查：我们只在这个``中间点'' \(\hat{x}\) 上计算梯度惩罚：
\end{enumerate}

\[
(\lVert \nabla D(\hat{x}) \rVert - 1)^2
\]

为什么这个``中间点''策略有效？想象 \(D\) 的``打分地貌''。我们强迫 \(D\) 在``真实国度''和``伪造国度''之间的整个广阔地带，都必须保持一个平滑的、坡度为 \(1\) 的``山坡''。这使得 \(G\) 无论身处``伪造国度''的哪个角落，都能清晰地看到一条通往``真实国度''的、坡度恒定的上山路径。而在这样的情况下，这个函数在``真实国度''与``伪造国度''的内部梯度极大概率也接近 \(1\)。

\Needspace{4\baselineskip}
\begin{enumerate}
\def\labelenumi{\arabic{enumi}.}
\setcounter{enumi}{2}
\tightlist
\item
  训练终点
\end{enumerate}

\Needspace{5\baselineskip}
我们之前在讨论 GANs 的目标函数时，推导出过这个结论：对于一个固定的 \(G\)，能让 \(V(D,G)\) 最大的那个``最优'' \(D^*\) 是：

\[
D^*(x)=\frac{p_{data}(x)}{p_{data}(x)+p_g(x)}
\]

\Needspace{5\baselineskip}
当 \(G\) 试图最小化 \(\max_D V(D,G)\) 时，它实际上是在最小化 \(C(G)\)：

\[
C(G)=2\cdot JSD(p_{data}\parallel p_g)-2\log 2
\]

所以，\(G\) 的目标（最小化 \(C(G)\)）等价于最小化 \(p_{data}\) 和 \(p_g\) 之间的 JSD 散度。

\Needspace{5\baselineskip}
当 \(p_{data}\) 和 \(p_g\) 完全不重叠时，JSD 会卡在恒定的最大值（\(\log 2\)），其梯度变为 \(0\)。证明：

\Needspace{5\baselineskip}
KL 散度 \(D_{KL}(P\parallel Q)\) 衡量从分布 \(P\) 切换到分布 \(Q\) 时``损失''了多少信息：

\[
D_{KL}(P\parallel Q)=\sum_x P(x)\log\left(\frac{P(x)}{Q(x)}\right)
\]

\Needspace{5\baselineskip}
JSD 散度只是 KL 散度的一个``平滑、对称''版本。首先定义一个``平均分布''：

\[
M=\frac{1}{2}(p_{data}+p_g)
\]

\Needspace{5\baselineskip}
然后：

\[
JSD(p_{data}\parallel p_g)=\frac{1}{2}D_{KL}(p_{data}\parallel M)+\frac{1}{2}D_{KL}(p_g\parallel M)
\]

\Needspace{5\baselineskip}
在``不重叠''的情况下，先看 \(D_{KL}(p_{data}\parallel M)\)：

\[
D_{KL}(p_{data}\parallel M)=\sum_x p_{data}(x)\log\left(\frac{p_{data}(x)}{M(x)}\right)
\]

\Needspace{5\baselineskip}
把 \(M(x)=\frac{1}{2}(p_{data}(x)+p_g(x))\) 代入：

\[
D_{KL}(p_{data}\parallel M)=\sum_x p_{data}(x)\log\left(\frac{p_{data}(x)}{\frac{1}{2}(p_{data}(x)+p_g(x))}\right)
\]

现在，我们只看那些 \(p_{data}(x)>0\) 的点；因为在其他地方 \(p_{data}(x)=0\)，整个求和项都为 0。

\Needspace{5\baselineskip}
在这些点上，由于 \(p_{data}(x)>0\) 且 \(p_g(x)=0\)，所以：

\[
\frac{p_{data}(x)}{\frac{1}{2}(p_{data}(x)+p_g(x))}=2
\]

\Needspace{5\baselineskip}
因此：

\[
D_{KL}(p_{data}\parallel M)=\sum_x p_{data}(x)\log 2=\log 2
\]

同理，\(D_{KL}(p_g\parallel M)=\log 2\)，所以 \(JSD(p_{data}\parallel p_g)=\log 2\)，变成了常数。故 \(C(G)\) 这时为常数，对 \(G\) 的梯度为 0，也就无法驱动生成器 \(G\) 的参数改变。

\Needspace{5\baselineskip}
\hypertarget{ux4e8cux566aux58f0ux6ce8ux5165}{%
\subsubsection{（二）噪声注入}\label{ux4e8cux566aux58f0ux6ce8ux5165}}

``噪声注入''不是一种单一、普适的模式崩溃解决方案，噪声加在何处决定了它的作用。向判别器输入添加 instance noise 可以平滑真实分布与生成分布，缓和二者支撑集几乎不重叠时的训练困难；在 StyleGAN 等模型的生成器中逐层注入噪声，主要用于控制头发、雀斑等随机细节，并不能据此断言它能普遍防止模式崩溃。因此，使用噪声时应明确具体方法及其目标，不能把``在生成器中间层加噪声''直接等同于通用的防模式崩溃机制。

\FloatBarrier
\Needspace{5\baselineskip}
\hypertarget{ux53c2ux8003ux6587ux732e-39}{%
\subsection{参考文献}\label{ux53c2ux8003ux6587ux732e-39}}

\vspace{0.25em}
\begin{itemize}
\tightlist
\item
  Goodfellow, I., Pouget-Abadie, J., Mirza, M., et al.~(2014). \href{https://papers.nips.cc/paper/5423-generative-adversarial-nets}{Generative Adversarial Nets}. NeurIPS 2014.
\item
  Arjovsky, M., \& Bottou, L. (2017). \href{https://arxiv.org/abs/1701.04862}{Towards Principled Methods for Training Generative Adversarial Networks}. ICLR 2017.
\item
  Arjovsky, M., Chintala, S., \& Bottou, L. (2017). \href{https://arxiv.org/abs/1701.07875}{Wasserstein GAN}. ICML 2017.
\item
  Gulrajani, I., Ahmed, F., Arjovsky, M., Dumoulin, V., \& Courville, A. (2017). \href{https://proceedings.neurips.cc/paper/2017/hash/892c3b1c6dccd52936e27cbd0ff683d6-Abstract.html}{Improved Training of Wasserstein GANs}. NeurIPS 2017.（WGAN-GP）
\item
  Karras, T., Laine, S., \& Aila, T. (2019). \href{https://openaccess.thecvf.com/content_CVPR_2019/html/Karras_A_Style-Based_Generator_Architecture_for_Generative_Adversarial_Networks_CVPR_2019_paper.html}{A Style-Based Generator Architecture for Generative Adversarial Networks}. CVPR 2019.
\end{itemize}

\clearpage
